\vspace{1.7mm}%  v4.3花形前距(0908实测回填)：课前预习花形为multicols栏首元素，行首胶被TeX弃置，前距须在栏区外给——设计4.5实出墨距2.82，+1.7 → 4.5±0.8窗
\begin{multicols}{2}
\emergencystretch=8em

\huaxing{课}{前}{预}{习}{知识导学\quad 素养初识}

\zsd{一}{空间向量的概念}

\tiaomu{1}{定义：在空间，我们把具有\kongda{大小}和\kongda{方向}的量叫做空间向量．\zhuzhu{与必修2平面向量初步的定义类比记忆(只多``空间''二字)；零向量方向任意、模为0，是概念题易错点。}}

\tiaomu{2}{(1)字母表示法：用字母\(\overrightarrow{a},\overrightarrow{b},\overrightarrow{c},\cdots\)表示．}

\bindp (2)几何表示法：用有向线段表示，其\kongda{长度}表示空间向量的模．即若向量\(\overrightarrow{a}\)的起点是\(A\)、终点是\(B\)，则向量\(\overrightarrow{a}\)也可记作\(\overrightarrow{AB}\)，其模记为\(| \overrightarrow{AB} |\)．\zhuzhu{字母表示与有向线段表示贯穿全章书写，解题前先统一记号；模的记号与绝对值区分，避免混淆。}

\tiaomu{3}{几类特殊向量(见下表)；规定：\kongda{零}向量与任意向量平行．即对任意向量\(\overrightarrow{a}\)，都有\(\overrightarrow{0}\parallel\overrightarrow{a}\)．\zhuzhu{概念辨析高频条目，判定按定义逐条核对；注意共线向量不要求在同一直线上(平行或重合均可)，且零向量与任意向量平行(衔接条目5)。}}

{\setlength{\tabcolsep}{2pt}\arrayrulecolor{gray122}
\vspace{1.1mm}
\noindent\begin{tabular}{!{\vline width 0.8pt}>{\fontsize{10.5pt}{13.9pt}\selectfont\centering\arraybackslash}m{16mm}!{\vline width 0.4pt}>{\fontsize{10.5pt}{13.9pt}\selectfont}m{35mm}!{\vline width 0.4pt}>{\fontsize{10.5pt}{13.9pt}\selectfont}m{26.7mm}!{\vline width 0.8pt}}
\thickhline
\makebox[\linewidth]{\textbf{名称}} & \makebox[\linewidth]{\textbf{定义}} & \makebox[\linewidth]{\textbf{表示}} \\[0.45mm]
\hline
零向量&长度为\kongda{0}的向量叫做零向量&记作\kongda{0} \\[0.45mm] \hline
单位向量&模等于\kongda{1}的向量&用\(e\)表示，\(|e|=1\) \\[0.45mm] \hline
相反向量&与\(\overrightarrow{a}\)长度相同而方向\kongda{相反}的向量&记作\(-\overrightarrow{a}\) \\[0.45mm] \hline
共线向量或平行向量&表示若干空间向量的有向线段所在的直线\kongda{互相平行}或\kongda{重合} &记作\(\overrightarrow{a}\parallel\overrightarrow{b}\) \\[0.45mm] \hline
相等向量&方向相同且\kongda{模相等}的向量&记作\(\overrightarrow{a}=\overrightarrow{b}\) \\[0.45mm]
\thickhline
\end{tabular}\par\addvspace{2.8mm}
}

\zhenhead{判断正误(正确的打\gou{},错误的打\cha{})}

\zhenti{(1)}{两个空间向量的模相等，则这两个向量相等．}{\cha{}}{模相等方向未必相同，还可能相反或斜交，两向量未必相等，故\cha{}．}

\zhenti{(2)}{相等向量一定是共线向量．}{\gou{}}{相等向量方向相同，所在直线平行或重合，符合共线向量(平行向量)定义，故\gou{}．}

\zsd{二}{空间向量的线性运算}

\tiaomu{1}{加法可按\kongda{三角形}法则或\kongda{平行四边形}法则作出；减法\(\overrightarrow{a}-\overrightarrow{b}\)为减向量终点指向被减向量终点；数乘\(\lambda\overrightarrow{a}\)：当\(\lambda>0\)时与\(\overrightarrow{a}\)方向\kongda{相同}，当\(\lambda<0\)时方向\kongda{相反}，当\(\lambda=0\)时\(\lambda\overrightarrow{a}=\overrightarrow{0}\)．\zhuzhu{全章运算的基础，加减、数乘的法则与运算律均与平面向量一致，类比迁移即可；线性组合的运用见条目9与10。}}

\tiaomu{2}{共线的充要条件：对任意两个空间向量\(\overrightarrow{a},\overrightarrow{b}\)(\(\overrightarrow{b}\neq\overrightarrow{0}\))，\(\overrightarrow{a}\parallel\overrightarrow{b}\)的充要条件是存在\kongda{实数} \(\lambda\)，使\(\overrightarrow{a}=\lambda\overrightarrow{b}\)．\zhuzhu{证线线平行与求参数存在性的基础工具；易错点：与条目9共面充要条件混淆------共线用一个实数、共面用有序实数对。}}

{\setlength{\tabcolsep}{3pt}\arrayrulecolor{gray122}
\vspace{1.1mm}
\noindent\begin{tabular}{!{\vline width 0.8pt}>{\fontsize{10.5pt}{13.9pt}\selectfont\centering\arraybackslash}m{10mm}!{\vline width 0.4pt}>{\fontsize{10.5pt}{13.9pt}\selectfont}m{37mm}!{\vline width 0.4pt}>{\fontsize{10.5pt}{13.9pt}\selectfont}m{28.6mm}!{\vline width 0.8pt}}
\thickhline
\makebox[\linewidth]{\textbf{运算}} & \makebox[\linewidth]{\textbf{法则要点}} & \makebox[\linewidth]{\textbf{运算律举例}} \\[0.45mm]
\hline
加法&三角形法则／平行四边形法则(与平面向量一致) & \(\overrightarrow{a}+\overrightarrow{b}=\overrightarrow{b}+\overrightarrow{a}\) \\[0.45mm] \hline
减法&减向量终点指向被减向量终点& \((\overrightarrow{a}+\overrightarrow{b})+\overrightarrow{c}=\overrightarrow{a}+(\overrightarrow{b}+\overrightarrow{c})\) \\[0.45mm] \hline
数乘&实数\(\lambda\)与向量\(\overrightarrow{a}\)的乘积仍为向量& \((\lambda+\mu)\overrightarrow{a}=\lambda\overrightarrow{a}+\mu\overrightarrow{a}\) \\[0.45mm] \hline
共线& \(\overrightarrow{a}\parallel\overrightarrow{b}\)(\(\overrightarrow{b}\neq\overrightarrow{0}\)) & \(\overrightarrow{a}=\lambda\overrightarrow{b}\) \\[0.45mm]
\thickhline
\end{tabular}\par\addvspace{2.8mm}
}

\zhenhead{判断正误(正确的打\gou{},错误的打\cha{})}

\zhenti{(1)}{若\(\overrightarrow{a}\parallel\overrightarrow{b}\)，\(\overrightarrow{b}\parallel\overrightarrow{c}\)，则\nobreak\(\overrightarrow{a}\parallel\overrightarrow{c}\)．}{\cha{}}{平行传递需\(\overrightarrow{b}\neq\overrightarrow{0}\)；\(\overrightarrow{b}=\overrightarrow{0}\)时零向量与任意向量平行，传递性失效，故\cha{}．}

\zhenti{(2)}{若\(\overrightarrow{a}\parallel\overrightarrow{b}\)，则存在唯一实数\(\lambda\)，使\(\overrightarrow{a}=\lambda\overrightarrow{b}\)．}{\cha{}}{缺\(\overrightarrow{b}\neq\overrightarrow{0}\)前提；且\(\overrightarrow{a}=\overrightarrow{0}\)时\(\lambda\)不唯一，故\cha{}．}

\zsd{三}{空间向量的夹角、数量积与共面}

\tiaomu{1}{夹角：已知两个非零向量\(\overrightarrow{a},\overrightarrow{b}\)，在空间中任取一点\(O\)，作\(\overrightarrow{OA}=\overrightarrow{a},\overrightarrow{OB}=\overrightarrow{b}\)，则大小在\kongda{[0°,180°]}内的\(\angle AOB\)称为\(\overrightarrow{a}\)与\(\overrightarrow{b}\)的夹角，记作\(\langle\overrightarrow{a},\overrightarrow{b}\rangle\)；当\(\langle\overrightarrow{a},\overrightarrow{b}\rangle=\frac{\pi}{2}\)时，称\(\overrightarrow{a}\)与\(\overrightarrow{b}\)\kongda{互相垂直}，记作\(\overrightarrow{a}\perp\overrightarrow{b}\)．\zhuzhu{夹角范围(0°到180°)是数量积正负讨论的前提；易混点：几何角(如异面直线所成角)与向量夹角可能互补，运算前先分清求的是哪个角。}}

\tiaomu{2}{数量积：定义\(\overrightarrow{a}\cdot\overrightarrow{b}=|\overrightarrow{a}||\overrightarrow{b}|\cos\langle\overrightarrow{a},\overrightarrow{b}\rangle\)；规定\kongda{零向量}与任意向量的数量积为0．性质：\(\overrightarrow{a}\perp\overrightarrow{b}\Leftrightarrow\overrightarrow{a}\cdot\overrightarrow{b}=0\)；\(\overrightarrow{a}\cdot\overrightarrow{a}={{\overrightarrow{a}}^{2}}={| \overrightarrow{a} |}^{2}\)；\(| \overrightarrow{a}\cdot\overrightarrow{b} |\leq| \overrightarrow{a} || \overrightarrow{b} |\)．\zhuzhu{零向量与任意向量的数量积为0．}\zhuzhu{求空间角与距离的核心运算(支撑条目16、17、34)；注意数量积不满足结合律，零向量与任意向量的数量积为0。}}

\tiaomu{3}{投影向量：(1)在空间，向量\(\overrightarrow{a}\)向向量\(\overrightarrow{b}\)投影：\zhuzhu{数量积几何意义的载体，条目33三余弦定理与条目38点面距的推导都基于投影；易错点：投影向量是向量、投影长度是数量。}}

\bindp 如图①，先将它们平移到同一平面\(\alpha\)内，利用平面上向量的投影，得到与向量\(\overrightarrow{b}\)共线的向量\(\overrightarrow{c}\)，\(\overrightarrow{c} =\)\(| \overrightarrow{a} |\cos\langle \overrightarrow{a},\overrightarrow{b} \rangle \cdot \frac{\overrightarrow{b}}{| \overrightarrow{b} |}\)，称向量\(\overrightarrow{c}\)为向量\(\overrightarrow{a}\)在向量\(\overrightarrow{b}\)上的投影向量．

\bindp (2)向量\(\overrightarrow{a}\)在直线\emph{l}上的投影如图②．

\bindp \par\vspace{1.9mm}\penalty10000\noindent\makebox[\linewidth][c]{\includegraphics[width=60mm,alt={@@@5da5e1ea-9d33-4579-a56f-804b4a72c894}]{media/media/sub3_B_4.png}}\par\vspace{-1.0mm}\penalty10000

\bindp (3)向量\(\overrightarrow{a}\)向平面\(\beta\)投影：

\bindp 如图③，分别由向量\(\overrightarrow{a}\)的起点\emph{A}和终点\emph{B}作平面\(\beta\)的垂线，垂足分别为\(A'\)，\(B'\)，得到向量\(\overrightarrow{A'B'}\)，向量\(\overrightarrow{A'B'}\)称为向量\(\overrightarrow{a}\)在平面\(\beta\)上的投影向量．

\tiaomu{4}{共面的充要条件：向量\(\overrightarrow{p}\)与不共线向量\(\overrightarrow{a},\overrightarrow{b}\)共面的充要条件是存在\kongda{有序实数对} \((x,y)\)，使\(\overrightarrow{p}=x\overrightarrow{a}+y\overrightarrow{b}\)．\zhuzhu{证四点共面、线共面的标准工具；易错点：定理前提是p与a、b不共线，此前提不可漏，且实数对存在唯一。}}

{\setlength{\tabcolsep}{3pt}\arrayrulecolor{gray122}
\vspace{1.1mm}
\noindent\begin{tabular}{!{\vline width 0.8pt}>{\fontsize{10.5pt}{13.9pt}\selectfont\centering\arraybackslash}m{10mm}!{\vline width 0.4pt}>{\fontsize{10.5pt}{13.9pt}\selectfont}m{36mm}!{\vline width 0.4pt}>{\fontsize{10.5pt}{13.9pt}\selectfont}m{29.6mm}!{\vline width 0.8pt}}
\thickhline
\makebox[\linewidth]{\textbf{概念}} & \makebox[\linewidth]{\textbf{要点}} & \makebox[\linewidth]{\textbf{记号／范围}} \\[0.45mm]
\hline
夹角& \(\angle AOB\)(\(\overrightarrow{OA}=\overrightarrow{a},\overrightarrow{OB}=\overrightarrow{b}\)) & \(0^\circ\leq\langle\overrightarrow{a},\overrightarrow{b}\rangle\leq180^\circ\) \\[0.45mm] \hline
数量积& \(\overrightarrow{a}\)的模与\(\overrightarrow{b}\)在\(\overrightarrow{a}\)上投影的数量之积& \(\overrightarrow{a}\allowbreak\cdot\allowbreak\overrightarrow{b}\allowbreak=\allowbreak|\overrightarrow{a}|\allowbreak|\overrightarrow{b}|\allowbreak\cos\allowbreak\langle\overrightarrow{a},\overrightarrow{b}\rangle\) \\[0.45mm] \hline
投影向量&与\(\overrightarrow{b}\)同向(\(\cos\allowbreak\langle\overrightarrow{a},\overrightarrow{b}\rangle\allowbreak>0\)时) &模为\(|\overrightarrow{a}|\allowbreak\cos\allowbreak\langle\overrightarrow{a},\overrightarrow{b}\rangle\) \\[0.45mm] \hline
共面&存在\((x,y)\)使\(\overrightarrow{p}=x\overrightarrow{a}+y\overrightarrow{b}\)(\(\overrightarrow{a},\overrightarrow{b}\)不共线) & \(\overrightarrow{p},\overrightarrow{a},\overrightarrow{b}\)共面\\[0.45mm]
\thickhline
\end{tabular}\par\addvspace{2.8mm}
}

\zhenhead{判断正误(正确的打\gou{},错误的打\cha{})}

\zhenti{(1)}{若\(\overrightarrow{p}=x\overrightarrow{a}+y\overrightarrow{b}\)(\(\overrightarrow{a}\)，\(\overrightarrow{b}\)不共线)，则\nobreak\(\overrightarrow{p}\)，\(\overrightarrow{a}\)，\(\overrightarrow{b}\)共面．}{\gou{}}{\(\overrightarrow{p}\)是\(\overrightarrow{a}\)、\(\overrightarrow{b}\)的线性组合，由共面向量定理知三向量共面，故\gou{}．}

\zhenti{(2)}{两个非零空间向量的数量积大于\(0\)，则这两个向量的夹角为锐角．}{\cha{}}{数量积大于0仅保证\(\cos\theta>0\)，夹角\(\theta\in[0^\circ,90^\circ)\)；\(\theta=0^\circ\)(同向)时数量积也为正而非锐角，故\cha{}．}

\huaxing{课}{中}{探}{究}{考点探究\quad 素养小结}

\tjdnr{一}{空间向量的概念辨析}{例\textbf{1}}{简单(知识点一)}{下列说法正确的是\nobreak(\kongwei)}

\bindopt {\fontsize{10.5pt}{19pt}\selectfont A．若\(| \overrightarrow{a} | = | \overrightarrow{b} |\)，则\nobreak\(\overrightarrow{a} = \overrightarrow{b}\)或\(\overrightarrow{a} = - \overrightarrow{b}\)；\par}

\bindopt {\fontsize{10.5pt}{19pt}\selectfont B．若\(\overrightarrow{a},\overrightarrow{b}\)为相反向量，则\nobreak\(\overrightarrow{a} + \overrightarrow{b} = \overrightarrow{0}\)；\par}

\bindopt {\fontsize{10.5pt}{19pt}\selectfont C．零向量是没有方向的向量；\par}

\bindopt {\fontsize{10.5pt}{19pt}\selectfont D．若\(\overrightarrow{a},\overrightarrow{b}\)是两个单位向量，则\nobreak\(\overrightarrow{a} = \overrightarrow{b}\)\par}

\bindp [分析]根据零向量、单位向量、向量的模、相等向量、相反向量等概念逐一分析即可得出答案.

\bindp [详解]解：若\(| \overrightarrow{a} | = | \overrightarrow{b} |\)，则它们的方向相同时是相等向量，方向相反时是相反向量，还有可能方向既不相同，也不相反，A错；

\bindp 若\(\overrightarrow{a},\overrightarrow{b}\)为相反向量，则它们的和为零向量，B对；

\bindp 零向量的方向是任意的，C错；

\bindp 两个单位向量只是模都为1，方向不一定相同，D错.故选：B.

\liB{变式\textbf{1}}{简单(知识点一)}{}{下列说法错误的是\nobreak(\kongwei)}

\bindopt {\fontsize{10.5pt}{19pt}\selectfont A．零向量与任意向量都平行；\par}

\bindopt {\fontsize{10.5pt}{19pt}\selectfont B．若\(|\overrightarrow{a}|=0\)，则\nobreak\(\overrightarrow{a}=\overrightarrow{0}\)；\par}

\bindopt {\fontsize{10.5pt}{19pt}\selectfont C．若\(\overrightarrow{a}=\overrightarrow{b}\)，则\nobreak\(|\overrightarrow{a}|=|\overrightarrow{b}|\)；\par}

\bindopt {\fontsize{10.5pt}{19pt}\selectfont D．若\(\overrightarrow{a}\parallel\overrightarrow{b}\)，则\nobreak\(|\overrightarrow{a}|=|\overrightarrow{b}|\)\par}

\ansline{\ansul{D}}

\jiexi{A规定成立；B模为0即零向量；C相等必等模；D平行不需等模(反例\(\overrightarrow{a}=2\overrightarrow{b}\))，故选D．}

\xiaojie{①识别：题干围绕模、方向、零向量、单位向量、相等向量与相反向量罗列说法要求辨误时，判归本题型。}

\bindp {\kaishu ②操作：对照定义逐条核验，存疑条目用反例否定。}

\bindp {\kaishu ③收束：排除错误项定答案，忌凭直觉误选。}

\tjdnr{二}{空间向量的线性运算}{例\textbf{1}}{简单(知识点二)}{在直三棱柱\(ABC - A_{1}B_{1}C_{1}\)中，若\(\overrightarrow{CA} = \overrightarrow{a},\overrightarrow{CB} = \overrightarrow{b},\overrightarrow{CC_{1}} = \overrightarrow{c},\)，则\nobreak\(\overrightarrow{A_{1}B}\)=\kongbai{}.(用\(\overrightarrow{a},\overrightarrow{b},\overrightarrow{c}\)表示)}

\bindp [分析]连接\(CA_{1},\)根据直三棱柱的结构特征及空间向量减法的几何意义可得\(\overrightarrow{A_{1}B} = \overrightarrow{CB} - \overrightarrow{CA} - \overrightarrow{CC_{1}}\)，结合已知即可求表达式.

\noindent\begin{minipage}[t]{49.000mm}\raggedright [详解]连接\(CA_{1},\)则\nobreak\(\overrightarrow{A_{1}B} = \overrightarrow{CB} - \overrightarrow{CA_{1}} = \overrightarrow{CB} - \overrightarrow{CA} - \overrightarrow{CC_{1}} =\overrightarrow{b} - \overrightarrow{a} - \overrightarrow{c}\).\end{minipage}\hspace{4.861mm}\begin{minipage}[t]{27.279mm}\centering\raisebox{\dimexpr-\height+3.329mm\relax}[\dimexpr3.329mm\relax][\dimexpr\height-3.329mm\relax]{\includegraphics[width=27.279mm,alt={@@@4e3aa752498e415b9cc3d6d60f7bb3fe}]{media/media/image1.png}}\end{minipage}\par

\bindp 故答案为：\(\overrightarrow{b} - \overrightarrow{a} - \overrightarrow{c}\)

\liB{变式\textbf{1}}{简单(知识点二)}{}{在平行六面体\(ABCD-A_1B_1C_1D_1\)中，\(\overrightarrow{AB}=\overrightarrow{a}\)，\(\overrightarrow{AD}=\overrightarrow{b}\)，\(\overrightarrow{AA_1}=\overrightarrow{c}\)，则\nobreak\(\overrightarrow{D_1B}=\)\kongbai{}．(用\(\overrightarrow{a}\)，\(\overrightarrow{b}\)，\(\overrightarrow{c}\)表示)}

\ansline{\ansul{\(\overrightarrow{a}-\overrightarrow{b}-\overrightarrow{c}\)}}

\jiexi{\(\overrightarrow{D_1B}=\overrightarrow{D_1A_1}+\overrightarrow{A_1A}+\overrightarrow{AB}=-\overrightarrow{b}-\overrightarrow{c}+\overrightarrow{a}\)．}

\xiaojie{①识别：在几何体中给出若干已知向量、要求表示指定向量时，判归本题型。}

\bindp {\kaishu ②操作：为目标向量选一条首尾相接的路径写出加减链，再把链中各段用相等或共线的已知向量替换，最后合并化简得表达式。}

\bindp {\kaishu ③收束：复查每段方向与起讫点，防抄错向量。}

\tjdnr{三}{用基底表示向量}{例\textbf{1}}{简单(知识点二)}{}

\noindent\begin{minipage}[t]{46.368mm}\raggedright 已知正方体\emph{ABCD}-\emph{A\textsubscript{1}B\textsubscript{1}C\textsubscript{1}D\textsubscript{1}}中，\(\overrightarrow{A_{1}E}\)=\(\frac{1}{4}\) \(\overrightarrow{A_{1}C_{1}}\)，若\(\overrightarrow{AE} =\)\emph{x}\(\overrightarrow{AA_{1}}\)+\emph{y}(\(\overrightarrow{AB}\)+\(\overrightarrow{AD}\))，则\emph{x}\nobreak=\kongbai{}，\emph{y}\nobreak=\kongbai{}.\end{minipage}\hspace{1.615mm}\begin{minipage}[t]{33.120mm}\centering\raisebox{\dimexpr-\height+3.557mm\relax}[\dimexpr3.557mm\relax][\dimexpr\height-3.557mm\relax]{\includegraphics[width=33.120mm,alt={@@@b0b55c9f-4634-4bba-b096-bc6dd0cf7747}]{media/media/image2.png}}\end{minipage}\par

\bindp [分析]结合空间向量的线性运算列方程，由此求得\(x,y\)的值.

\bindp [详解]\(\overrightarrow{AE} = \overrightarrow{AA_{1}} + \overrightarrow{A_{1}E} = \overrightarrow{AA_{1}} + \frac{1}{4}\overrightarrow{A_{1}C_{1}} = \overrightarrow{AA_{1}} + \frac{1}{4}( \overrightarrow{AB} + \overrightarrow{AD} )\)，所以\(x = 1,y = \frac{1}{4}\).故答案为：\(1\)；\(\frac{1}{4}\)

\liB{变式\textbf{1}}{简单(知识点二)}{}{在平行六面体\(ABCD-A_1B_1C_1D_1\)中，点\(E\)在\(A_1C_1\)上，且\(\overrightarrow{A_1E}=\frac{1}{2}\overrightarrow{A_1C_1}\)，若\(\overrightarrow{AE}=x\overrightarrow{AA_1}+y(\overrightarrow{AB}+\overrightarrow{AD})\)，则\nobreak\(x=\)\kongbai{}，\(y=\)\kongbai{}．}

\ansline{\ansul{\(x=1\)}，\ansul{\(y=\frac{1}{2}\)}}

\jiexi{\(\overrightarrow{AE}=\overrightarrow{AA_1}+\overrightarrow{A_1E}=\overrightarrow{AA_1}+\frac{1}{2}(\overrightarrow{AB}+\overrightarrow{AD})\)，故\(x=1\)，\(y=\frac{1}{2}\)．}

\xiaojie{①识别：指定三个不共面的向量作基底、要求把目标向量写成基底向量的线性组合并求系数时，判归本题型。}

\bindp {\kaishu ②操作：先把目标向量沿几何路径分解为基底向量的和，再把各段按比例关系代入基底。}

\bindp {\kaishu ③收束：比较等式两边系数求出待定值，防漏写中间段。}

\tjdnr{四}{共面向量的判定}{例\textbf{1}}{简单(知识点三)}{有下列说法:}

\bindopt ①若\(\overrightarrow{p} = x\overrightarrow{a} + y\overrightarrow{b}\)，则\nobreak\(\overrightarrow{p}\)与\(\overrightarrow{a}\)，\(\overrightarrow{b}\)共面；②若\(\overrightarrow{p}\)与\(\overrightarrow{a}\)，\(\overrightarrow{b}\)共面，则\nobreak\(\overrightarrow{p}\)=\emph{x}\(\overrightarrow{a}\)+\emph{y}\(\overrightarrow{b}\)；③若\(\overrightarrow{MP}\)=\emph{x}\(\overrightarrow{MA}\)+\emph{y}\(\overrightarrow{MB}\)，则\emph{P}，\emph{M}，\emph{A}，\emph{B}共面；④若\emph{P}，\emph{M}，\emph{A}，\emph{B}共面，则\nobreak\(\overrightarrow{MP}\)=\emph{x}\(\overrightarrow{MA}\)+\emph{y}\(\overrightarrow{MB}\).其中正确的是\nobreak(\kongwei)

\bindopt {\fontsize{10.5pt}{19pt}\selectfont \makebox[\dimexpr0.5\linewidth-1em\relax][l]{A．①②③④；}\makebox[\dimexpr0.5\linewidth-1em\relax][l]{B．①③④；}\par}

\bindopt {\fontsize{10.5pt}{19pt}\selectfont \makebox[\dimexpr0.5\linewidth-1em\relax][l]{C．①③；}\makebox[\dimexpr0.5\linewidth-1em\relax][l]{D．②④}\par}

\bindp [分析]利用空间向量共面定理逐一判断即可.

\bindp [详解]若\(\overrightarrow{a}\)，\(\overrightarrow{b}\)共线，由\(\overrightarrow{p}\)=\emph{x}\(\overrightarrow{a}\)+\emph{y}\(\overrightarrow{b}\)知\(\overrightarrow{p}\)一定与\(\overrightarrow{a}\)，\(\overrightarrow{b}\)共面，

\bindp 若\(\overrightarrow{a}\)，\(\overrightarrow{b}\)不共线，则满足共面定理，\(\overrightarrow{p}\)与\(\overrightarrow{a}\)，\(\overrightarrow{b}\)共面，①对；

\bindp 同理③对；若\(\overrightarrow{p}\)与\(\overrightarrow{a}\)，\(\overrightarrow{b}\)共面，且\(\overrightarrow{a}\)，\(\overrightarrow{b}\)共线，则不一定有\(\overrightarrow{p}\)=\emph{x}\(\overrightarrow{a}\)+\emph{y}\(\overrightarrow{b}\)，故②不对；同理④不对，故选：C.

\liB{变式\textbf{1}}{简单(知识点三)}{}{已知\(\overrightarrow{MP}=3\overrightarrow{MA}-2\overrightarrow{MB}\)，判断\(P\)，\(M\)，\(A\)，\(B\)四点是否共面，并说明理由．}

\ansline{\ansul{共面}(\(\overrightarrow{MP}\)可由\(\overrightarrow{MA}\)，\(\overrightarrow{MB}\)线性表示，故\(\overrightarrow{MP}\)，\(\overrightarrow{MA}\)，\(\overrightarrow{MB}\)共面，即\(P\)在平面\(MAB\)内)}

\jiexi{\(\overrightarrow{MP}=3\overrightarrow{MA}-2\overrightarrow{MB}\)为\(\overrightarrow{MA}\)、\(\overrightarrow{MB}\)的线性组合，故三向量共面，\(P\)在平面\(MAB\)内．}

\xiaojie{①识别：题干给出p=xa+yb或MP=xMA+yMB一类线性表示与“共面”结论的互推说法判断时，判归本题型。}

\bindp {\kaishu ②操作：核对充要条件的前提(两向量不共线或三点不共线)是否齐备，再对缺前提的说法举共线反例排除。}

\bindp {\kaishu ③收束：最后汇总正确序号定选项。}

\tjdnr{五}{数量积的概念与运算律}{例\textbf{1}\duoxuan}{简单(知识点三)}{设\(\overrightarrow{a}\)，\(\overrightarrow{b}\)为空间中的任意两个非零向量，下列各式中正确的有(\kongwei)}

\bindopt {\fontsize{10.5pt}{19pt}\selectfont A．\({\overrightarrow{a}}^{2} = | \overrightarrow{a} |^{2}\)；\par}

\bindopt {\fontsize{10.5pt}{19pt}\selectfont B．\(\frac{\overrightarrow{a} \cdot \overrightarrow{b}}{\overrightarrow{a} \cdot \overrightarrow{a}} = \frac{\overrightarrow{b}}{\overrightarrow{a}}\)；\par}

\bindopt {\fontsize{10.5pt}{19pt}\selectfont C．\(( \overrightarrow{a} \cdot \overrightarrow{b} )^{2} = {\overrightarrow{a}}^{2} \cdot {\overrightarrow{b}}^{2}\)；\par}

\bindopt {\fontsize{10.5pt}{19pt}\selectfont D．\(( \overrightarrow{a} - \overrightarrow{b} )^{2} = {\overrightarrow{a}}^{2} - 2\overrightarrow{a} \cdot \overrightarrow{b} + {\overrightarrow{b}}^{2}\)\par}

\bindp [分析]根据空间向量数量积的定义与运算律一一判断即可；

\bindp [详解]解：对于A：\({\overrightarrow{a}}^{2} = \overrightarrow{a} \cdot \overrightarrow{a} = | \overrightarrow{a} | \cdot | \overrightarrow{a} |\cos 0 = | \overrightarrow{a} |^{2}\)，故A正确；

\bindp 对于B：因为向量不能做除法，即\(\frac{\overrightarrow{b}}{\overrightarrow{a}}\)无意义，故B错误；

\bindp 对于C：\(( \overrightarrow{a} \cdot \overrightarrow{b} )^{2} = ( | \overrightarrow{a} | \cdot | \overrightarrow{b} |\cos\langle \overrightarrow{a},\overrightarrow{b} \rangle )^{2} = | \overrightarrow{a} |^{2} \cdot | \overrightarrow{b} |^{2}\cos^{2}\langle \overrightarrow{a},\overrightarrow{b} \rangle\)，故C错误；

\bindp 对于D：\(( \overrightarrow{a} - \overrightarrow{b} )^{2} = ( \overrightarrow{a} - \overrightarrow{b} ) \cdot ( \overrightarrow{a} - \overrightarrow{b} ) = {\overrightarrow{a}}^{2} - 2\overrightarrow{a} \cdot \overrightarrow{b} + {\overrightarrow{b}}^{2}\)，故D正确；故选：AD

\liB{变式\textbf{1}}{简单(知识点三)}{}{已知\(|\overrightarrow{a}|=2\)，\(|\overrightarrow{b}|=3\)，\(\langle\overrightarrow{a},\overrightarrow{b}\rangle=60^\circ\)，则\nobreak\((\overrightarrow{a}+2\overrightarrow{b})\cdot(\overrightarrow{a}-\overrightarrow{b})=\)\kongbai{}．}

\ansline{\ansul{\(-11\)}}

\jiexi{\(\overrightarrow{a}\cdot\overrightarrow{b}=2\times3\times\cos60^\circ=3\)；原式\(=\overrightarrow{a}^2+\overrightarrow{a}\cdot\overrightarrow{b}-2\overrightarrow{b}^2=4+3-18=-11\)．}

\xiaojie{①识别：题干列出数量积相关的若干等式(平方、商式、乘积展开)要求判断正误时，判归本题型。}

\bindp {\kaishu ②操作：按定义a·b=|a||b|cosθ逐式展开核验，展开后移项验算。}

\bindp {\kaishu ③收束：警惕向量无除法、数量积无结合律等误用。}

\tjdnr{六}{数量积求夹角与投影}{例\textbf{1}}{中档(知识点三)}{}

\noindent\begin{minipage}[t]{46.368mm}\raggedright 如图，在正方体\(ABCD - A_{1}B_{1}C_{1}D_{1}\)中，求向量\(\overrightarrow{BC_{1}}\)与\(\overrightarrow{AC}\)的夹角的大小．\end{minipage}\hspace{2.502mm}\begin{minipage}[t]{33.120mm}\centering\raisebox{\dimexpr-\height+3.719mm\relax}[\dimexpr3.719mm\relax][\dimexpr\height-3.719mm\relax]{\includegraphics[width=33.120mm,alt={@@@d8ffc038-cfd6-4d16-b846-f3d81f29b325}]{media/media/image3.png}}\end{minipage}\par

\bindp [分析]方法1：结合图形，平移向量，利用空间向量的夹角定义来求，但要注意向量夹角的范围；

\bindp 方法2：先求\(\overrightarrow{BC_{1}} \cdot \overrightarrow{AC}\)，再利用公式\(\cos\langle \overrightarrow{BC_{1}},\overrightarrow{AC} \rangle = \frac{\overrightarrow{BC_{1}} \cdot \overrightarrow{AC}}{| \overrightarrow{BC_{1}} | \cdot | \overrightarrow{AC} |}\)求\(\cos\langle \overrightarrow{BC_{1}},\overrightarrow{AC} \rangle\)，最后确定\(\langle \overrightarrow{BC_{1}},\overrightarrow{AC} \rangle\)即可．

\bindp [详解]解：方法1：因为\(\overrightarrow{AD_{1}} = \overrightarrow{BC_{1}}\)，所以\(\angle CAD_{1}\)的大小就等于\(\langle \overrightarrow{BC_{1}},\overrightarrow{AC} \rangle\)

\bindp 因为△\(CAD_{1}\)为等边三角形，所以\(\angle CAD_{1} = 60^{{^\circ}}\)，所以\(\overrightarrow{BC_{1}}\)与\(\overrightarrow{AC}\)的夹角的大小为\(60{^\circ}\)．

\bindp 方法2．设正方体的棱长为1，\(\overrightarrow{BC_{1}} \cdot \overrightarrow{AC} = ( \overrightarrow{BC} + \overrightarrow{CC_{1}} ) \cdot ( \overrightarrow{AB} + \overrightarrow{BC} ) = ( \overrightarrow{AD} + \overrightarrow{AA_{1}} ) \cdot ( \overrightarrow{AB} + \overrightarrow{AD} )\)\(= \overrightarrow{AD} \cdot \overrightarrow{AB} + | \overrightarrow{AD} |^{2} + \overrightarrow{AA_{1}} \cdot \overrightarrow{AB} + \overrightarrow{AA_{1}} \cdot \overrightarrow{AD} = 0 + | \overrightarrow{AD} |^{2} + 0 + 0 = | \overrightarrow{AD} |^{2} = 1\)又因为\(| \overrightarrow{BC_{1}} | = \sqrt{2},| \overrightarrow{AC} | = \sqrt{2}\)，所以\(\cos\langle \overrightarrow{BC_{1}},\overrightarrow{AC} \rangle = \frac{\overrightarrow{BC_{1}} \cdot \overrightarrow{AC}}{| \overrightarrow{BC_{1}} | \cdot | \overrightarrow{AC} |} = \frac{1}{\sqrt{2} \times \sqrt{2}} = \frac{1}{2}\)，
因为\(\langle \overrightarrow{BC_{1}},\overrightarrow{AC} \rangle \in \lbrack 0,\pi\rbrack\)，所以\(\overrightarrow{BC_{1}}\)与\(\overrightarrow{AC}\)的夹角的大小为\(60{^\circ}\)．

\liB{变式\textbf{1}}{中档(知识点三)}{}{在长方体\(ABCD-A_1B_1C_1D_1\)中，\(AB=2\)，\(AD=2\)，\(AA_1=1\)，则向量\(\overrightarrow{AB_1}\)与\(\overrightarrow{AC}\)夹角的余弦值为\kongbai{}．}

\ansline{\ansul{\(\frac{\sqrt{10}}{5}\)}}

\jiexi{\(\overrightarrow{AB_1}\cdot\overrightarrow{AC}=(\overrightarrow{AB}+\overrightarrow{BB_1})\cdot(\overrightarrow{AB}+\overrightarrow{BC})=|\overrightarrow{AB}|^2=4\)，\(|\overrightarrow{AB_1}|=\sqrt{5}\)，\(|\overrightarrow{AC}|=2\sqrt{2}\)，故\(\cos\langle\overrightarrow{AB_1},\overrightarrow{AC}\rangle=\frac{4}{\sqrt{5}\times2\sqrt{2}}=\frac{\sqrt{10}}{5}\)．}

\xiaojie{①识别：题干在几何体中求两向量的夹角、模长或投影而不建坐标系时，判归本题型。}

\bindp {\kaishu ②操作：先平移向量使起点共点或选不共面的三个向量作基底，再把相关向量用基底表示并算出数量积与模长。}

\bindp {\kaishu ③收束：最后代入夹角公式并按向量夹角范围[0,π]取角。}

\tjdnr{七}{数量积条件求参}{例\textbf{1}}{中档(知识点三)}{已知向量\(\overrightarrow{a},\overrightarrow{b}\)，且\(| \overrightarrow{a} | = 2\)，\(| \overrightarrow{b} | = 1\)，\(\langle \overrightarrow{a},\overrightarrow{b} \rangle = 60^{\circ}\)，则使向量\(\overrightarrow{a} + \lambda\overrightarrow{b}\)与\(\lambda\overrightarrow{a} - 2\overrightarrow{b}\)的夹角为钝角的实数\(\lambda\)的取值范围是\kongbai{}．}

\bindp [分析]根据\(( \overrightarrow{a} + \lambda\overrightarrow{b} ) \cdot ( \lambda\overrightarrow{a} - 2\overrightarrow{b} ) < 0\)，得出方程求得\(- 1 - \sqrt{3} < \lambda < - 1 + \sqrt{3}\)，若向量\(\overrightarrow{a} + \lambda\overrightarrow{b}\)与\(\lambda\overrightarrow{a} - 2\overrightarrow{b}\)反向共线时，则存在实数\(k < 0\)使得\(\overrightarrow{a} + \lambda\overrightarrow{b} = k(\lambda\overrightarrow{a} - 2\overrightarrow{b})\)成立，得到方程组\(\Bigg\{ \begin{array}{r}
k\lambda = 1 \\
\lambda = - 2k
\end{array} \Bigg.\)无解，即可得到答案.

\bindp [详解]由向量\(\overrightarrow{a},\overrightarrow{b}\)，且\(| \overrightarrow{a} | = 2\)，\(| \overrightarrow{b} | = 1\)，\(\langle \overrightarrow{a},\overrightarrow{b} \rangle = 60^{\circ}\)，可得\(\overrightarrow{a} \cdot \overrightarrow{b} = 2 \times 1 \times \cos 60^{\circ} = 1\)，

\bindp 因为向量\(\overrightarrow{a} + \lambda\overrightarrow{b}\)与\(\lambda\overrightarrow{a} - 2\overrightarrow{b}\)的夹角为钝角，可得\(\cos\langle \overrightarrow{a} + \lambda\overrightarrow{b},\lambda\overrightarrow{a} - 2\overrightarrow{b} \rangle < 0\)且\(\cos\langle \overrightarrow{a} + \lambda\overrightarrow{b},\lambda\overrightarrow{a} - 2\overrightarrow{b} \rangle \neq - 1\)，由\(( \overrightarrow{a} + \lambda\overrightarrow{b} ) \cdot ( \lambda\overrightarrow{a} - 2\overrightarrow{b} ) < 0\)，即\(\lambda{\overrightarrow{a}}^{2}\text{+(}\lambda^{2} - 2)\overrightarrow{a} \cdot \overrightarrow{b} - 2\lambda{\overrightarrow{b}}^{2} < 0\)，所以\(\lambda^{2} + 2\lambda - 2 < 0\)，解得\(- 1 - \sqrt{3} < \lambda < - 1 + \sqrt{3}\)，

\bindp 若向量\(\overrightarrow{a} + \lambda\overrightarrow{b}\)与\(\lambda\overrightarrow{a} - 2\overrightarrow{b}\)反向共线时，存在实数\(k < 0\)，使得\(\overrightarrow{a} + \lambda\overrightarrow{b} = k(\lambda\overrightarrow{a} - 2\overrightarrow{b})\)成立，可得\(\Bigg\{ \begin{array}{r}
k\lambda = 1 \\
\lambda = - 2k
\end{array} \Bigg.\)，此时方程组无解，所以不存在实数\(k\)使得向量\(\overrightarrow{a} + \lambda\overrightarrow{b}\)与\(\lambda\overrightarrow{a} - 2\overrightarrow{b}\)反向共线，所以实数\(\lambda\)的取值范围是\(( - 1 - \sqrt{3}, - 1 + \sqrt{3} )\).

\liB{变式\textbf{1}}{简单(知识点三)}{}{已知\(| \overrightarrow{a} | = 3\sqrt{2},| \overrightarrow{b} | = 4,\overrightarrow{m} = \overrightarrow{a} + \overrightarrow{b},\overrightarrow{n} = \overrightarrow{a} + \lambda\overrightarrow{b},\langle \overrightarrow{a},\overrightarrow{b} \rangle = 135{^\circ},\overrightarrow{m}\bot\overrightarrow{n}\)，则λ=\kongbai{}．}

\ansline{\ansul{\(- \frac{3}{2}\)}}

\jiexi{\(\overrightarrow{m}\cdot\overrightarrow{n}=(\overrightarrow{a}+\overrightarrow{b})\cdot(\overrightarrow{a}+\lambda\overrightarrow{b})=|\overrightarrow{a}|^2+\lambda|\overrightarrow{b}|^2+(1+\lambda)\overrightarrow{a}\cdot\overrightarrow{b}=18+16\lambda-12(1+\lambda)=6+4\lambda=0\)，故\(\lambda=-\frac{3}{2}\)．}

\xiaojie{①识别：题干给出两向量的模与夹角、并以夹角为锐角/钝角或两向量垂直为条件求参数时，判归本题型。}

\bindp {\kaishu ②操作：第一步用数量积把条件翻译成不等式或方程(锐角cos>0、钝角cos<0、垂直cos=0)。}

\bindp {\kaishu ③收束：再求解并剔除使两向量共线的参数，最后写出参数值或取值范围。}

\tjdnr{八}{数量积求距离(展开法)}{例\textbf{1}}{中档(知识点三)}{}

\noindent\begin{minipage}[t]{44.000mm}\raggedright 如图，在大小为45°的二面角A\-EF\-D中，四边形ABFE，CDEF都是边长为1的正方形，则B，D两点间的距离是\nobreak\mbox{}\nobreak\hfill(\kongwei)\hspace{0.56mm}\end{minipage}\hspace{5.055mm}\begin{minipage}[t]{33.120mm}\centering\raisebox{\dimexpr-\height+4.093mm\relax}[\dimexpr4.093mm\relax][\dimexpr\height-4.093mm\relax]{\includegraphics[width=33.120mm,alt={@@@494a628b97d74c2894c135f69e8f1eb7}]{media/media/image4.png}}\end{minipage}\par

\bindopt {\fontsize{10.5pt}{19pt}\selectfont \makebox[\dimexpr0.5\linewidth-1em\relax][l]{A．\(\sqrt{3}\)；}\makebox[\dimexpr0.5\linewidth-1em\relax][l]{B．\(\sqrt{2}\)；}\par}

\bindopt {\fontsize{10.5pt}{19pt}\selectfont \makebox[\dimexpr0.5\linewidth-1em\relax][l]{C．1；}\makebox[\dimexpr0.5\linewidth-1em\relax][l]{D．\(\sqrt{3 - \sqrt{2}}\)}\par}

\bindp [分析]由\(\overrightarrow{DB} = \overrightarrow{ED} + \overrightarrow{FE} + \overrightarrow{BF}\)，利用数量积运算性质展开即可得到答案

\bindp [详解]\(\because\overrightarrow{BD} = \overrightarrow{ED} + \overrightarrow{FE} + \overrightarrow{BF}\),\(\therefore| \overrightarrow{BD} |^{2} = | \overrightarrow{BF} |^{2} + | \overrightarrow{FE} |^{2} + | \overrightarrow{ED} |^{2} + 2\overrightarrow{BF} \bullet \overrightarrow{FE} + 2\overrightarrow{FE} \bullet \overrightarrow{ED} + 2\overrightarrow{BF} \bullet \overrightarrow{ED} = 1 + 1 + 1 - \sqrt{2}\)故\(| \overrightarrow{BD} | = \sqrt{3 - \sqrt{2}}\)故选\(D\)

\bindp [点睛]本题是要求空间两点之间的距离，运用空间向量将其表示，然后计算得到结果，较为基础．

\liB{变式\textbf{1}}{简单(知识点三)}{}{在大小为\(60^\circ\)的二面角\(A-EF-D\)中，四边形\(ABFE\)，\(CDEF\)都是边长为\(1\)的正方形，则\nobreak\(B\)，\(D\)两点间的距离为\kongbai{}．}

\ansline{\ansul{\(\sqrt{2}\)}}

\jiexi{\(|\overrightarrow{BD}|^2=1+1+1+2\overrightarrow{BF}\cdot\overrightarrow{ED}=3+2\cos120^\circ=2\)，故\(|\overrightarrow{BD}|=\sqrt{2}\)．}

\xiaojie{①识别：题干在含二面角的拼接图形中求两点距离时，判归本题型。}

\bindp {\kaishu ②操作：第一步把待求线段对应的向量写成已知棱向量的首尾相接和链，再对模长平方作数量积展开，逐项代入模长与夹角(端点处两条垂线段的夹角由二面角决定，注意取45°还是135°)。}

\bindp {\kaishu ③收束：最后开方得距离。}

\tjdnr{九}{数量积求距离(折叠矩形)}{例\textbf{1}}{中档(知识点三)}{已知矩形\(ABCD\)中，\(AB = 1\)，\(BC = \sqrt{3}\)，将矩形\(ABCD\)沿对角线\(AC\)折起，使平面\(ABC\)与平面\(ACD\)垂直，则\nobreak\(|\overrightarrow{BD}| =\)(\kongwei)．}

\bindopt {\fontsize{10.5pt}{19pt}\selectfont \makebox[\dimexpr0.25\linewidth-0.5em\relax][l]{A．\(\frac{\sqrt{10}}{2}\)；}\makebox[\dimexpr0.25\linewidth-0.5em\relax][l]{B．\(\frac{\sqrt{6}}{2}\)；}\makebox[\dimexpr0.25\linewidth-0.5em\relax][l]{C．\(\frac{\sqrt{5}}{2}\)；}\makebox[\dimexpr0.25\linewidth-0.5em\relax][l]{D．2}\par}

\bindp [分析]过点\(B\)，\(D\)分别向\(AC\)作垂线，垂足分别为\(M\)，\(N\)，由\(\overrightarrow{BD} = \overrightarrow{BM} + \overrightarrow{MN} + \overrightarrow{ND}\)，平方后结合长度和垂直关系可得解.

\noindent\begin{minipage}[t]{46.368mm}\raggedright [详解]过点\(B\)，\(D\)分别向\(AC\)作垂线，垂足分别为\(M\)，\(N\)，则可得\(AM = \frac{1}{2}\)，\(BM = \frac{\sqrt{3}}{2}\)，\(CN = \frac{1}{2}\)，\(DN = \frac{\sqrt{3}}{2}\)，\(MN = 1\)．由于\(\overrightarrow{BD} = \overrightarrow{BM} + \overrightarrow{MN} + \overrightarrow{ND}\)，所以\(|\overrightarrow{BD}|^{2} = {(\overrightarrow{BM} + \overrightarrow{MN} + \overrightarrow{ND})}^{2}\)\(= |\overrightarrow{BM}|^{2} + |\overrightarrow{MN}|^{2} + |\overrightarrow{ND}|^{2} + 2(\overrightarrow{BM} \cdot \overrightarrow{MN} + \overrightarrow{MN} \cdot \overrightarrow{ND} + \overrightarrow{BM} \cdot \overrightarrow{ND})\)\(= ( \frac{\sqrt{3}}{2} )^{2} + 1^{2} + ( \frac{\sqrt{3}}{2} )^{2} + 2 \times (0 + 0 + 0) = \frac{5}{2}\)，所以\(|\overrightarrow{BD}| = \frac{\sqrt{10}}{2}\)．故选:\emph{A}．\end{minipage}\hspace{3.832mm}\begin{minipage}[t]{30.000mm}\centering\raisebox{\dimexpr-\height+3.677mm\relax}[\dimexpr3.677mm\relax][\dimexpr\height-3.677mm\relax]{\includegraphics[width=30.000mm,alt={@@@e34e5148-bca3-440e-84b1-3cbf33fe2cc1}]{media/media/image5.png}}\end{minipage}\par

\bindp [点睛]本题主要考查了空间向量数量积的应用，属于基础题.

\liB{变式\textbf{1}}{中档(知识点三)}{}{已知矩形\(ABCD\)中，\(AB=1\)，\(BC=\sqrt{3}\)，将矩形\(ABCD\)沿对角线\(AC\)折起，使二面角\(B-AC-D\)的大小为\(60^\circ\)，则\nobreak\(|\overrightarrow{BD}|=\)\kongbai{}．}

\ansline{\ansul{\(\frac{\sqrt{13}}{2}\)}}

\jiexi{\(BM=DN=\frac{\sqrt{3}}{2}\)，\(MN=1\)；两垂线向量夹角随二面角为\(60^\circ\)，\(\overrightarrow{BM}\cdot\overrightarrow{ND}=\frac{3}{4}\cos60^\circ=\frac{3}{8}\)；\(|\overrightarrow{BD}|^2=\frac{3}{4}+1+\frac{3}{4}+2\times\frac{3}{8}=\frac{13}{4}\)，故\(|\overrightarrow{BD}|=\frac{\sqrt{13}}{2}\)．}

\xiaojie{①识别：题干为平面图形沿一条直线折起且折后两面垂直、求两点间距离时，判归本题型。}

\bindp {\kaishu ②操作：第一步在折后图形中过两点对折痕作垂线定出垂足，再把目标向量写成两段垂线与折痕段的和链并平方展开。}

\bindp {\kaishu ③收束：最后利用面面垂直带来的向量垂直逐项化简、开方得距离。}

\vbox{

\huaxing{课}{堂}{评}{价}{知识评价\quad 素养形成}

\jiancestem{{\fontsize{11.4pt}{14pt}\selectfont\heihao {\numboldjian 1}．}\hspace{2.7mm}\tieside{简单(知识点三)}\hspace{2.2mm}已知\(|\overrightarrow{a}|=2\)，\(|\overrightarrow{b}|=3\)，\(\langle\overrightarrow{a},\overrightarrow{b}\rangle=60^\circ\)，则\nobreak\(\overrightarrow{a}\cdot\overrightarrow{b}=\)\kongbai{}．{\zhushi{(提示：\(\overrightarrow{a}\cdot\overrightarrow{b}=|\overrightarrow{a}||\overrightarrow{b}|\cos\langle\overrightarrow{a},\overrightarrow{b}\rangle\))}}}

\ansline{\ansul{3}}

\jiexi{\(\overrightarrow{a}\cdot\overrightarrow{b}=|\overrightarrow{a}||\overrightarrow{b}|\cos60^\circ=2\times3\times0.5=3\)．}

\jiancestem{{\fontsize{11.4pt}{14pt}\selectfont\heihao {\numboldjian 2}．}\hspace{2.7mm}\tieside{简单(知识点二)}\hspace{2.2mm}设\(\overrightarrow{e_1}\)，\(\overrightarrow{e_2}\)不共线，\(\overrightarrow{a}=2\overrightarrow{e_1}+\overrightarrow{e_2}\)，\(\overrightarrow{b}=4\overrightarrow{e_1}+k\overrightarrow{e_2}\)，若\(\overrightarrow{a}\parallel\overrightarrow{b}\)，则\nobreak\(k=\)(\kongwei)}

\bindopt {\fontsize{10.5pt}{21pt}\selectfont \makebox[\dimexpr0.25\linewidth-0.5em\relax][l]{A．\(\frac{1}{2}\)；}\makebox[\dimexpr0.25\linewidth-0.5em\relax][l]{B．\(2\)；}\makebox[\dimexpr0.25\linewidth-0.5em\relax][l]{C．\(-\frac{1}{2}\)；}\makebox[\dimexpr0.25\linewidth-0.5em\relax][l]{D．\(3\)}\par}

\ansline{\ansul{B}}

\jiexi{\(\overrightarrow{b}=\lambda\overrightarrow{a}\Rightarrow4=2\lambda\Rightarrow\lambda=2\)，\(k=2\)，选B．}

\jiancestem{{\fontsize{11.4pt}{14pt}\selectfont\heihao {\numboldjian 3}．}\hspace{2.7mm}\tieside{简单(知识点三)}\hspace{2.2mm}设\(\overrightarrow{a}\)，\(\overrightarrow{b}\)均为非零空间向量，则“\(\overrightarrow{a}\cdot\overrightarrow{b}=0\)”是“\(\overrightarrow{a}\perp\overrightarrow{b}\)”的(\kongwei)}

\bindopt {\fontsize{10.5pt}{21pt}\selectfont A．充分不必要条件；\par}

\bindopt {\fontsize{10.5pt}{21pt}\selectfont B．必要不充分条件；\par}

\bindopt {\fontsize{10.5pt}{21pt}\selectfont C．充要条件；\par}

\bindopt {\fontsize{10.5pt}{21pt}\selectfont D．既不充分也不必要条件\par}

\ansline{\ansul{C}}

\jiexi{非零向量\(\overrightarrow{a}\cdot\overrightarrow{b}=0\Leftrightarrow\overrightarrow{a}\perp\overrightarrow{b}\)，充要条件，选C．}

\jiancestem{{\fontsize{11.4pt}{14pt}\selectfont\heihao {\numboldjian 4}．}\hspace{2.7mm}\tieside{简单(知识点二)}\hspace{2.2mm}在空间四边形\(ABCD\)中，\(E\)，\(F\)分别为\(AB\)，\(CD\)的中点，则\nobreak\(\frac{1}{2}(\overrightarrow{AD}+\overrightarrow{BC})=\)(\kongwei)}

\bindopt {\fontsize{10.5pt}{21pt}\selectfont \makebox[\dimexpr0.25\linewidth-0.5em\relax][l]{A．\(2\overrightarrow{EF}\)；}\makebox[\dimexpr0.25\linewidth-0.5em\relax][l]{B．\(-\overrightarrow{EF}\)；}\makebox[\dimexpr0.25\linewidth-0.5em\relax][l]{C．\(\overrightarrow{FE}\)；}\makebox[\dimexpr0.25\linewidth-0.5em\relax][l]{D．\(\overrightarrow{EF}\)}\par}

\ansline{\ansul{D}}

\jiexi{取\(AC\)中点\(G\)，\(\overrightarrow{EF}=\overrightarrow{EG}+\overrightarrow{GF}=\frac{1}{2}\overrightarrow{BC}+\frac{1}{2}\overrightarrow{AD}\)，选D．}

\jiance{{\fontsize{11.4pt}{14pt}\selectfont\heihao {\numboldjian 5}．}\hspace{2.7mm}\tieside{简单(知识点三)}\hspace{2.2mm}已知\(|\overrightarrow{a}|=4\)，\(|\overrightarrow{b}|=3\)，\(\langle\overrightarrow{a},\overrightarrow{b}\rangle=120^\circ\)，则\nobreak\(|\overrightarrow{a}+\overrightarrow{b}|=\)\kongbai{}．{\zhushi{(提示：先算\(\overrightarrow{a}\cdot\overrightarrow{b}\)，再对\(|\overrightarrow{a}+\overrightarrow{b}|^2=(\overrightarrow{a}+\overrightarrow{b})^2\)展开)}}}{\ansul{\(\sqrt{13}\)}}

\jiexi{\(\overrightarrow{a}\cdot\overrightarrow{b}=4\times3\times\cos120^\circ=-6\)，\(|\overrightarrow{a}+\overrightarrow{b}|^2=16+9-12=13\)，故\(\sqrt{13}\)．}

\vspace{-9.7pt}% 字替对照 C：课堂评价 vbox 在 10.09/18.25 密度下自然高超出栏高 9.58pt，于盒尾胶内收缩（不动任何字号/行距/表距）
}

\end{multicols}
